\documentclass[11pt]{article}
\usepackage{amsmath, amssymb, amsthm}
\usepackage{geometry}
\usepackage{hyperref}
\usepackage{bm}

\geometry{margin=1in}

\title{Anima-Vectorium v4.0: Runtime-Aligned Mathematical Backbone for the Spectral Core}
\author{Borealis S. Hedling}
\date{October 2026}

\begin{document}
\maketitle

\begin{abstract}
This document formalizes the mathematical backbone of the Anima-Vectorium
architecture in its v4.0 form. Unlike v3.0 Mathématique, which provided a
conceptual geometric foundation, v4.0 introduces runtime-aligned operator
definitions, a computational Laplacian suitable for mobile thermodynamic
constraints, membrane invariants compatible with reversible operator algebra,
and a bisimulation framework linking expressive and spectral states. This
version is designed explicitly to support Option C (Rust Core + Mobile
Bindings), ensuring deterministic execution, bounded expressive dynamics, and
thermodynamically safe operator injection. Academic references are included to
support future expressive-clarity versions.
\end{abstract}

\section{Introduction}

Anima-Vectorium models expressive cognition as a geometric dynamical system
evolving on a manifold $M$ equipped with a connection $\nabla$, a Laplacian
operator $\Delta$, and a set of reversible operators $\mathcal{O}$. The v4.0
architecture aligns these mathematical objects with the constraints of a
mobile runtime, including thermodynamic envelopes, bounded drift, and
membrane-governed operator injection.

This version introduces:
\begin{itemize}
    \item A runtime-aligned Laplacian $\Delta_R$ suitable for SIMD execution.
    \item A 9-dimensional expressive manifold $M = \mathbb{R}^9$ with bounded
    norm.
    \item A reversible operator algebra $\mathcal{O}$ with membrane invariants.
    \item A drift clamp defined via rotational stability constraints.
    \item A bisimulation relation linking expressive and spectral states.
    \item A thermodynamic envelope defined via power budget and spectral variance.
\end{itemize}

\section{Expressive Manifold}

We define the expressive manifold as:
\[
M = \mathbb{R}^9
\]
with expressive state:
\[
\bm{e} \in M, \qquad \|\bm{e}\|_2 \leq 1.
\]

The bounded norm ensures:
\begin{itemize}
    \item stability under mobile thermodynamic constraints,
    \item compatibility with glyph projection,
    \item reversible operator application.
\end{itemize}

\section{Runtime-Aligned Laplacian}

We define the runtime Laplacian $\Delta_R$ as:
\[
\Delta_R f = \sum_{i=1}^9 w_i \frac{\partial^2 f}{\partial x_i^2}
\]
where $w_i$ are thermodynamic weights satisfying:
\[
0 < w_i \leq 1, \qquad \sum_i w_i = 1.
\]

This form is chosen because:
\begin{itemize}
    \item it is SIMD-friendly,
    \item it is compatible with mobile power envelopes,
    \item it preserves spectral stability.
\end{itemize}

\section{Operator Algebra}

Define operator algebra:
\[
\mathcal{O} = \{ O : M \rightarrow M \mid O \text{ reversible},\; O \text{ bounded} \}.
\]

Each operator $O$ satisfies:
\[
\|O(\bm{e})\|_2 \leq 1,
\]
and reversibility:
\[
O^{-1} \in \mathcal{O}.
\]

Operators are injected via packets:
\[
O(\bm{e}) = \bm{e} + \alpha \bm{d},
\]
where $\alpha$ is magnitude and $\bm{d}$ is direction.

\section{Membrane Invariants}

Define membrane state:
\[
m \in \{0,1,2,3\}
\]
corresponding to neutral, active, clamp, reset.

Membrane invariants:
\[
m = 2 \Rightarrow O(\bm{e}) = \bm{e}
\]
\[
m = 3 \Rightarrow \bm{e} = \bm{e}_0
\]

Clamp threshold:
\[
\alpha > \alpha_{\max} \Rightarrow m = 2.
\]

Reset threshold:
\[
\|\bm{e}\|_2 > 1 \Rightarrow m = 3.
\]

\section{Drift Clamp}

Define drift vector:
\[
\bm{d}_{\text{drift}} = \nabla \bm{e}.
\]

Clamp condition:
\[
\|\bm{d}_{\text{drift}}\|_2 > \delta_{\max} \Rightarrow m = 2.
\]

Rotational stability:
\[
\theta = \arccos\left( \frac{\bm{e} \cdot \bm{d}_{\text{drift}}}{\|\bm{e}\|\|\bm{d}_{\text{drift}}\|} \right)
\]
\[
\theta > \theta_{\max} \Rightarrow m = 2.
\]

\section{Thermodynamic Envelope}

Define power envelope:
\[
P = \sum_{i=1}^9 w_i \left( \frac{\partial^2 f}{\partial x_i^2} \right)^2.
\]

Envelope condition:
\[
P \leq P_{\max}.
\]

Spectral variance:
\[
\sigma^2 = \sum_{i=1}^9 \left( \frac{\partial f}{\partial x_i} \right)^2.
\]

Envelope violation:
\[
\sigma^2 > \sigma_{\max}^2 \Rightarrow m = 2.
\]

\section{Bisimulation}

Define bisimulation relation:
\[
\bm{e} \sim \bm{s}
\]
if:
\[
\Delta_R \bm{e} = \Delta_R \bm{s}
\]
and
\[
m(\bm{e}) = m(\bm{s}).
\]

This ensures expressive and spectral states evolve consistently.

\section{Conclusion}

v4.0 Mathématique provides a runtime-aligned mathematical backbone suitable for
Option C. It formalizes the Laplacian, operator algebra, membrane invariants,
drift clamps, thermodynamic envelope, and bisimulation in a manner compatible
with deterministic Rust execution and mobile constraints.

\section{References}

\begin{thebibliography}{9}

\bibitem{doCarmo}
M. P. do Carmo,
\textit{Riemannian Geometry},
Birkhäuser, 1992.

\bibitem{leeSmooth}
J. M. Lee,
\textit{Introduction to Smooth Manifolds},
Springer, 2012.

\bibitem{evansPDE}
L. C. Evans,
\textit{Partial Differential Equations},
American Mathematical Society, 2010.

\bibitem{rudinFA}
W. Rudin,
\textit{Functional Analysis},
McGraw-Hill, 1991.

\bibitem{hairerODE}
E. Hairer, S. P. Nørsett, G. Wanner,
\textit{Solving Ordinary Differential Equations I},
Springer, 1993.

\bibitem{sastryNonlinear}
S. Sastry,
\textit{Nonlinear Systems: Analysis, Stability, and Control},
Springer, 1999.

\bibitem{boydConvex}
S. Boyd, L. Vandenberghe,
\textit{Convex Optimization},
Cambridge University Press, 2004.

\end{thebibliography}

\end{document}
